% Slide thuyết trình bảo vệ Niên luận / Khóa luận tốt nghiệp
% Đề tài: Nghiên cứu kỹ thuật khai thác và phòng chống tấn công Man-in-the-Middle
% Sinh viên: Nguyễn Thị Kỳ Duyên (MSSV: 0023410647 - Lớp: ĐHCNTT23B-CS)
% GVHD: TS. Lương Thái Ngọc
% Build:  .\build.ps1 slide-bao-ve.tex (chạy từ thư mục gốc)

\documentclass[fleqn]{beamer}
\usepackage[utf8]{inputenc}
\usepackage[T5]{fontenc}
\usepackage[vietnamese]{babel}
\usepackage{mathtools,amssymb,amsmath}
\usepackage{graphicx}
\usepackage{booktabs}
\usepackage{tabularx}
\newcolumntype{L}{>{\raggedright\arraybackslash}X}
\usepackage{tikz}
\usetikzlibrary{shapes,arrows.meta,positioning}
\usepackage[backend=biber,style=numeric]{biblatex}
\addbibresource{tai-lieu-tham-khao.bib}

\hypersetup{unicode}

\usetheme{Madrid}
\usecolortheme{default}
\graphicspath{{images/}}

\AtBeginSection[]{\begin{frame}\frametitle{Nội dung}\tableofcontents[currentsection]\end{frame}}

\title[Nghiên cứu Tấn công \& Phòng chống MitM]{NGHIÊN CỨU KỸ THUẬT KHAI THÁC VÀ PHÒNG CHỐNG TẤN CÔNG MAN-IN-THE-MIDDLE}
\author[Nguyễn Thị Kỳ Duyên]{Sinh viên thực hiện: \textbf{Nguyễn Thị Kỳ Duyên} \\ \small MSSV: 0023410647 -- Lớp: ĐHCNTT23B-CS}
\institute[DTHU]{Khoa Công nghệ và Kỹ thuật, Trường Đại học Đồng Tháp\\ Ngành: Khoa học máy tính -- Chuyên ngành: Mạng Máy Tính và An Ninh\\ \textbf{Giảng viên hướng dẫn: TS. Lương Thái Ngọc}}
\date{Năm học 2026 -- 2027}

\begin{document}

\begin{frame}\titlepage\end{frame}

\begin{frame}{Nội dung trình bày}
  \tableofcontents
\end{frame}

\section{Đặt vấn đề \& Mục tiêu nghiên cứu}

\begin{frame}{Tính cấp thiết của đề tài}
  \begin{itemize}
    \item \textbf{Bối cảnh an toàn mạng hiện đại:}
      \begin{itemize}
        \item Mọi giao dịch ngân hàng, định danh cá nhân và dữ liệu doanh nghiệp đều luân chuyển qua hạ tầng mạng máy tính.
        \item Môi trường mạng LAN truyền thống thường bị ngộ nhận là ``an toàn nội bộ''.
      \end{itemize}
    \item \textbf{Điểm yếu cố hữu của các giao thức TCP/IP sơ khai:}
      \begin{itemize}
        \item Giao thức ARP (RFC 826), DNS (RFC 1035), DHCP (RFC 2131) \textbf{thiếu hoàn toàn cơ chế xác thực nguồn gốc}.
        \item Các thiết bị mù quáng tiếp nhận thông tin phản hồi mà không kiểm tra tính xác thực.
      \end{itemize}
    \item \textbf{Nguy cơ từ tấn công Man-in-the-Middle (MitM):}
      \begin{itemize}
        \item Kẻ tấn công bí mật chèn vào giữa kênh truyền, nghe lén mật khẩu, trích xuất Cookie, chiếm đoạt phiên và giả mạo nội dung.
      \end{itemize}
  \end{itemize}
\end{frame}

\begin{frame}{Mục tiêu và phạm vi đề tài}
  \begin{columns}
    \begin{column}{0.5\textwidth}
      \begin{block}{Mục tiêu nghiên cứu}
        \begin{itemize}
          \item Hệ thống hóa các kỹ thuật dẫn đường: ARP, DNS, DHCP Spoofing, SSL Stripping.
          \item Phân tích các kỹ thuật tấn công phiên: Session Hijacking, Cookie Theft, Tampering.
          \item Dựng Lab thực nghiệm 5 node và triển khai tấn công kiểm chứng.
          \item Cấu hình giải pháp phòng thủ: HTTPS/HSTS, DAI, DHCP Snooping, VPN, Arpwatch.
        \end{itemize}
      \end{block}
    \end{column}
    \begin{column}{0.5\textwidth}
      \begin{block}{Phạm vi đề tài}
        \begin{itemize}
          \item Môi trường thử nghiệm an toàn độc lập (Isolated Testbed Lab).
          \item Thiết bị: Kali Linux, 02 máy nạn nhân (Win/Linux), Gateway, Web Server và Switch quản lý.
          \item Đo kiểm định lượng trước và sau khi kích hoạt các giải pháp an ninh.
        \end{itemize}
      \end{block}
    \end{column}
  \end{columns}
\end{frame}

\section{Các kỹ thuật khai thác \& dẫn đường MitM}

\begin{frame}{Khái niệm tấn công Man-in-the-Middle}
  \begin{center}
    \begin{tikzpicture}[scale=0.8, every node/.style={transform shape}]
      \node[draw, rectangle, fill=blue!15, minimum width=2.2cm, minimum height=1cm, rounded corners] (vic) at (0,0) {Victim};
      \node[draw, rectangle, fill=red!25, minimum width=2.6cm, minimum height=1.2cm, rounded corners, line width=1.2pt, align=center] (atk) at (4.5,0) {\textbf{Attacker (MitM)}\\\footnotesize(Sniffing / Tampering)};
      \node[draw, rectangle, fill=green!15, minimum width=2.2cm, minimum height=1cm, rounded corners] (gw) at (9,0) {Gateway / Server};

      \draw[<->, >=Stealth, line width=1.5pt, red] (vic) -- node[above]{\footnotesize Đánh cắp tin} (atk);
      \draw[<->, >=Stealth, line width=1.5pt, red] (atk) -- node[above]{\footnotesize Chuyển tiếp} (gw);
    \end{tikzpicture}
  \end{center}
  \begin{itemize}
    \item \textbf{Bản chất:} Kẻ tấn công đứng giữa luồng truyền thông, giả mạo Gateway đối với Victim và giả mạo Victim đối với Gateway.
    \item \textbf{Mục tiêu:} Nghe lén (Eavesdropping), đánh cắp định danh (Credentials), chiếm quyền phiên (Session Hijacking), giả mạo nội dung (Tampering).
  \end{itemize}
\end{frame}

\begin{frame}{Kỹ thuật ARP Spoofing / Poisoning}
  \begin{itemize}
    \item \textbf{Nguyên lý:} Gửi liên tục các gói tin \textbf{ARP Reply giả mạo hai chiều}:
      \begin{itemize}
        \item Nói với Victim: \textit{``IP Gateway hiện đang mang địa chỉ MAC của Attacker''}.
        \item Nói với Gateway: \textit{``IP Victim hiện đang mang địa chỉ MAC của Attacker''}.
      \end{itemize}
    \item \textbf{Kết quả:} Bảng ARP Cache của cả hai bên bị nhiễm độc.
    \item \textbf{Bật chuyển tiếp gói tin (IP Forwarding):}
      \texttt{sysctl -w net.ipv4.ip\_forward=1} $\rightarrow$ Lưu lượng đi qua máy Attacker một cách trong suốt, không làm rớt mạng nạn nhân.
  \end{itemize}
\end{frame}

\begin{frame}{Kỹ thuật DNS Spoofing, DHCP Spoofing \& SSL Stripping}
  \begin{itemize}
    \item \textbf{DNS Spoofing:} Đón đầu truy vấn UDP cổng 53, trả lời bản ghi A Record giả mạo trỏ về Web Server lừa đảo của Attacker.
    \item \textbf{DHCP Spoofing (Rogue DHCP):} Dựng máy chủ DHCP lậu, phản hồi nhanh hơn máy chủ thật để cấp phát Default Gateway và DNS Server trỏ thẳng về Attacker.
    \item \textbf{SSL Stripping (Moxie Marlinspike):}
      \begin{itemize}
        \item Đón đầu gói tin chuyển hướng HTTP 301/302.
        \item Duy trì kết nối HTTPS với Server thật nhưng hạ cấp xuống \textbf{HTTP bản rõ} khi phản hồi cho Victim.
        \item Nạn nhân không thấy ổ khóa HTTPS nhưng trang web vẫn hiển thị, mọi mật khẩu nhập vào đều bị lộ.
      \end{itemize}
  \end{itemize}
\end{frame}

\begin{frame}{Các kiểu tấn công vào phiên và dữ liệu}
  \begin{columns}
    \begin{column}{0.5\textwidth}
      \begin{block}{Session Hijacking \& Cookie Theft}
        \begin{itemize}
          \item Bắt gói tin HTTP chứa Cookie phiên (\texttt{PHPSESSID}, Token).
          \item Tiêm Cookie vào trình duyệt máy Attacker để truy cập thẳng hệ thống không cần mật khẩu.
          \item Khai thác Cookie thiếu cờ \texttt{Secure} và \texttt{HttpOnly}.
        \end{itemize}
      \end{block}
    \end{column}
    \begin{column}{0.5\textwidth}
      \begin{block}{Giả mạo nội dung (Content Tampering)}
        \begin{itemize}
          \item Chèn thẻ mã độc JavaScript (\texttt{<script src="hook.js">}) kết nối trình duyệt vào khung BeEF.
          \item Sửa đổi tham số giao dịch trong form POST (thay đổi số tài khoản nhận tiền).
        \end{itemize}
      \end{block}
    \end{column}
  \end{columns}
\end{frame}

\section{Mô hình phòng thử nghiệm (Testbed Lab)}

\begin{frame}{Kiến trúc mô hình Testbed Lab 5 thành phần}
  \begin{center}
    \begin{tikzpicture}[scale=0.75, every node/.style={transform shape}]
      \node[draw, rectangle, fill=blue!20, minimum width=8cm, minimum height=1cm, rounded corners] (sw) at (5,2.5) {\textbf{Cisco Catalyst Switch (VLAN 10)} --- Bật DHCP Snooping \& DAI};

      \node[draw, rectangle, fill=green!15, minimum width=2.5cm, minimum height=1cm, align=center] (gw) at (0,4.5) {\textbf{Gateway (VyOS)}\\\footnotesize 192.168.10.1};
      \node[draw, rectangle, fill=green!15, minimum width=2.5cm, minimum height=1cm, align=center] (web) at (10,4.5) {\textbf{Web Server}\\\footnotesize 192.168.10.100};

      \node[draw, rectangle, fill=blue!10, minimum width=2.4cm, minimum height=1cm, align=center] (v1) at (0,0.5) {\textbf{Victim 1 (Win)}\\\footnotesize 192.168.10.10};
      \node[draw, rectangle, fill=red!20, minimum width=2.6cm, minimum height=1cm, align=center] (at) at (5,0.5) {\textbf{Kali Linux}\\\footnotesize 192.168.10.50};
      \node[draw, rectangle, fill=blue!10, minimum width=2.4cm, minimum height=1cm, align=center] (v2) at (10,0.5) {\textbf{Victim 2 (Linux)}\\\footnotesize 192.168.10.20};

      \draw (gw) -- (sw); \draw (web) -- (sw);
      \draw (v1) -- (sw); \draw (at) -- (sw); \draw (v2) -- (sw);
    \end{tikzpicture}
  \end{center}
  \begin{itemize}
    \item \textbf{Phân đoạn mạng:} \texttt{192.168.10.0/24} (VLAN 10 an toàn độc lập).
    \item \textbf{Cổng Switch:} Gi1/0/1 (Gateway - \textbf{Trusted}), Fa1/0/1 - Fa1/0/3 (\textbf{Untrusted}).
  \end{itemize}
\end{frame}

\section{Giải pháp phòng chống \& Kiểm chứng thực nghiệm}

\begin{frame}{Giải pháp 1: HTTPS kết hợp cơ chế HSTS}
  \begin{itemize}
    \item \textbf{Cấu hình Nginx Web Server:}
      \begin{itemize}
        \item Kích hoạt giao thức TLS 1.3 và chứng chỉ số X.509 hợp lệ.
        \item Bổ sung tiêu đề HSTS bắt buộc:
      \end{itemize}
      \texttt{Strict-Transport-Security: max-age=31536000; includeSubDomains; preload}
    \item \textbf{Kiểm chứng chạy lại tấn công SSL Stripping:}
      \begin{itemize}
        \item Trình duyệt của Victim tự động nâng cấp nội bộ lên HTTPS (\texttt{Internal 307 Redirect}) trước khi gửi ra mạng.
        \item Kẻ tấn công Bettercap không thể đón đầu yêu cầu HTTP.
        \item \textbf{Kết quả:} Tấn công SSL Stripping bị \textbf{vô hiệu hóa 100\%}.
      \end{itemize}
  \end{itemize}
\end{frame}

\begin{frame}{Giải pháp 2: DHCP Snooping trên Switch}
  \begin{itemize}
    \item \textbf{Cấu hình trên Cisco Catalyst Switch:}
      \begin{itemize}
        \item \texttt{ip dhcp snooping} và \texttt{ip dhcp snooping vlan 10}.
        \item Cổng kết nối Gateway hợp pháp: \texttt{ip dhcp snooping trust}.
        \item Cổng kết nối người dùng/Attacker: Mặc định là \textbf{Untrusted}.
      \end{itemize}
    \item \textbf{Kiểm chứng chạy lại tấn công Rogue DHCP Server:}
      \begin{itemize}
        \item Kali phát tán gói tin DHCP Offer giả mạo trên cổng Untrusted.
        \item Switch phát hiện và lập tức \textbf{DROP toàn bộ gói tin Offer} từ cổng này.
        \item \textbf{Kết quả:} Nạn nhân 100\% chỉ nhận IP từ Gateway thật; ngăn chặn triệt để Rogue DHCP.
      \end{itemize}
  \end{itemize}
\end{frame}

\begin{frame}{Giải pháp 3: Dynamic ARP Inspection (DAI)}
  \begin{itemize}
    \item \textbf{Cơ chế hoạt động:}
      \begin{itemize}
        \item Switch đối chiếu từng gói ARP trên cổng Untrusted với \textbf{DHCP Snooping Binding Table}.
        \item Nếu cặp IP-MAC không khớp: Lập tức \textbf{DROP} và ghi log cảnh báo vi phạm.
      \end{itemize}
    \item \textbf{Kiểm chứng chạy lại tấn công ARP Cache Poisoning:}
      \begin{itemize}
        \item Kali chạy kịch bản \texttt{arpspoof} mạo nhận IP Gateway \texttt{192.168.10.1}.
        \item Switch phát hiện MAC Kali không trùng khớp với IP Gateway trong bảng Binding.
        \item Switch ghi nhận log: \texttt{\%SW\_DAI-4-DHCP\_SNOOPING\_DENY: 1 invalid ARPs...}
        \item \textbf{Kết quả:} Bảng ARP của Victim hoàn toàn trong sạch; ARP Spoofing bị triệt tiêu tận gốc.
      \end{itemize}
  \end{itemize}
\end{frame}

\begin{frame}{Giải pháp 4 \& 5: VPN WireGuard \& Giám sát ARP bất thường}
  \begin{columns}
    \begin{column}{0.5\textwidth}
      \begin{block}{Đường hầm mã hóa VPN}
        \begin{itemize}
          \item Kích hoạt kết nối VPN WireGuard giữa Victim và Gateway.
          \item Toàn bộ dữ liệu được mã hóa bằng thuật toán ChaCha20-Poly1305.
          \item Bắt gói tin trên mạng LAN chỉ thu được gói tin UDP mã hóa, không thể đọc dữ liệu.
        \end{itemize}
      \end{block}
    \end{column}
    \begin{column}{0.5\textwidth}
      \begin{block}{Giám sát ARP bất thường}
        \begin{itemize}
          \item Vận hành Arpwatch và kịch bản phân tích ARP thời gian thực.
          \item Phát hiện ngay hiện tượng xung đột MAC (MAC Flapping) của Gateway.
          \item Gửi cảnh báo tức thì tới quản trị viên để cách ly thiết bị vi phạm.
        \end{itemize}
      \end{block}
    \end{column}
  \end{columns}
\end{frame}

\section{Đánh giá kết quả \& Kết luận}

\begin{frame}{Bảng tổng hợp kết quả đối sánh định lượng}
  \begin{table}[htbp]
    \centering
    \footnotesize
    \begin{tabularx}{\textwidth}{l L c c}
      \toprule
      \textbf{Kịch bản kiểm thử} & \textbf{Chỉ số đo lường} & \textbf{Trước phòng vệ} & \textbf{Sau khi bảo vệ} \\
      \midrule
      \textbf{ARP Poisoning} & Tỷ lệ làm sai lệch bảng ARP & 100\% (Nhiễm độc) & \textbf{0\% (Bị DAI Drop)} \\
      \textbf{Lộ mật khẩu HTTP} & Số tài khoản bị lộ bản rõ & 10/10 lần & \textbf{0/10 (Mã hóa TLS)} \\
      \textbf{Rogue DHCP} & Nạn nhân nhận IP lậu & 90\% & \textbf{0\% (Bị Drop)} \\
      \textbf{SSL Stripping} & Hạ cấp thành công kết nối & 10/10 lần & \textbf{0/10 (HSTS chặn)} \\
      \textbf{Cookie / Session} & Chiếm quyền phiên thành công & 100\% & \textbf{0\% (Bảo vệ)} \\
      \textbf{Sniffing qua VPN} & Tỷ lệ dữ liệu bị đọc lén & 100\% & \textbf{0\% (UDP mã hóa)} \\
      \bottomrule
    \end{tabularx}
  \end{table}
\end{frame}

\begin{frame}{Kết luận và Đóng góp của đề tài}
  \begin{itemize}
    \item \textbf{Kết luận:}
      \begin{itemize}
        \item Tấn công Man-in-the-Middle là mối đe dọa nghiêm trọng khai thác trực diện điểm yếu không xác thực của các giao thức mạng cổ điển.
        \item Không có một giải pháp đơn lẻ nào đủ an toàn; bắt buộc phải triển khai mô hình \textbf{phòng thủ theo chiều sâu (Defense-in-Depth)}.
      \end{itemize}
    \item \textbf{Đóng góp của đề tài:}
      \begin{itemize}
        \item Hệ thống hóa toàn diện lý thuyết tấn công và phòng chống MitM.
        \item Xây dựng thành công mô hình Testbed Lab 5 node kiểm chứng trực quan.
        \item Cung cấp bộ kịch bản PoC và cấu hình mẫu sẵn sàng ứng dụng cho hệ thống mạng trường học và doanh nghiệp.
      \end{itemize}
  \end{itemize}
\end{frame}

\begin{frame}{Hướng phát triển tiếp theo}
  \begin{enumerate}
    \item Nghiên cứu mở rộng các kỹ thuật MitM trên giao thức thế hệ mới \textbf{IPv6} (như Rogue Router Advertisement, Neighbor Discovery Spoofing) và cơ chế RA Guard.
    \item Ứng dụng \textbf{Trí tuệ nhân tạo (Machine Learning)} để phát hiện bất thường luồng lưu lượng MitM theo thời gian thực.
    \item Tích hợp cơ chế phòng vệ tự động vào kiến trúc mạng điều khiển bằng phần mềm (\textbf{SDN}).
  \end{enumerate}
\end{frame}

\begin{frame}{Đáp ứng tiêu chuẩn đánh giá Niên luận (Mẫu Khoa 2026)}
  \begin{columns}
    \begin{column}{0.34\textwidth}
      \begin{block}{\textbf{1. Quyển báo cáo (4.0 đ)}}
        \footnotesize
        \begin{itemize}
          \item \textbf{Hình thức (1.0 đ):} Đúng chuẩn lề 3-2-2-2, Times 13pt, giãn dòng 1.3, trích dẫn IEEE/BibLaTeX 100\% nhất quán.
          \item \textbf{Nội dung (3.0 đ):} 4 chương mạch lạc, phân tích sâu 6 kỹ thuật dẫn đường và 5 giải pháp phòng vệ.
        \end{itemize}
      \end{block}
    \end{column}
    \begin{column}{0.34\textwidth}
      \begin{block}{\textbf{2. Sản phẩm (4.0 đ)}}
        \footnotesize
        \begin{itemize}
          \item \textbf{Giao diện (1.0 đ):} Hệ sinh thái Docsify trực quan, lab web testbed DVWA thân thiện.
          \item \textbf{Demo mô hình (3.0 đ):} Testbed 5 node, PoC chạy mượt mà, cấu hình Cisco DAI, HSTS, VPN và script cảnh báo ARP.
        \end{itemize}
      \end{block}
    \end{column}
    \begin{column}{0.32\textwidth}
      \begin{block}{\textbf{3. Trình bày (2.0 đ)}}
        \footnotesize
        \begin{itemize}
          \item \textbf{Thuyết trình:} Mạch lạc, tự tin, slide Beamer trực quan.
          \item \textbf{Phản biện:} Nắm vững bản chất gói tin tầng 2 và tầng 7, trả lời thuyết phục các câu hỏi.
        \end{itemize}
      \end{block}
    \end{column}
  \end{columns}
\end{frame}

\begin{frame}
  \centering
  \vspace{1cm}
  {\LARGE\bfseries XIN TRÂN TRỌNG CẢM ƠN}\\[0.5cm]
  {\large Quý Thầy, Cô trong Hội đồng và các bạn đã chú ý lắng nghe!}\\[1.5cm]
  {\large\textit{Em rất mong nhận được những câu hỏi và ý kiến đóng góp của Quý Thầy, Cô!}}
\end{frame}

\end{document}
